% !TEX program = xelatex
\documentclass[12pt,a4paper]{ctexart}

\usepackage{fontspec}
\setCJKmainfont{SimSun}[BoldFont=SimHei,ItalicFont=KaiTi]
\setCJKsansfont{SimHei}
\setCJKmonofont{FangSong}

\usepackage[top=28mm,bottom=28mm,left=28mm,right=28mm]{geometry}
\usepackage{fancyhdr}
\usepackage{perpage}
\MakePerPage{footnote}

\setlength{\parindent}{2em}
\setlength{\parskip}{0pt}

\newcommand{\circledfootnotenumber}{%
  {\fontspec{SimSun}\ifcase\value{footnote}\or ①\or ②\or ③\or ④\or ⑤\or ⑥\or ⑦\or ⑧\or ⑨\else\arabic{footnote}\fi}}
\renewcommand{\thefootnote}{\circledfootnotenumber}

\pagestyle{fancy}
\fancyhf{}
\fancyfoot[C]{\thepage}
\renewcommand{\headrulewidth}{0pt}

\begin{document}

\begin{center}
    {\zihao{-3}\heiti 半个世纪后，我想对您说：从认识中国到建设中国\par}
    \vspace{0.8em}
    {\zihao{-4}\kaishu 姓名：\underline{\hspace{3.5cm}}\qquad 学号：\underline{\hspace{4cm}}\par}
\end{center}

\vspace{1em}

{\zihao{-4}\songti\setlength{\baselineskip}{20pt}\selectfont

毛泽东同志：

2026年是您逝世五十周年。写下“毛泽东同志”这几个字时，我其实有一点停顿：这不像平时写读书报告，直接分析观点就可以了；它更像是要把这一学期学到的东西，放到一个具体的倾诉对象面前重新想一遍。刚开始上这门课时，我对“毛泽东思想和中国特色社会主义理论体系概论”的理解比较简单，觉得它大概就是教材概念、课堂笔记和期末作业。现在回头看，这种想法有些浅了。真正让我留下印象的，不是某一个需要背下来的定义，而是课程一直在追问的那个问题：中国为什么会走上自己的道路，又怎样在变化的时代里继续走下去。

我这学期印象最深的一次阅读，是《新民主主义论》。当时选择这篇文章，一部分原因是它在课程中出现得多，另一部分原因是它开头的“中国向何处去”很直接。这个问题放在1940年的中国，并不是一句修辞，而是民族危机压到眼前时必须回答的问题。读第一遍时，我更多是在圈出“半殖民地半封建社会”“两步走”“新政治、新经济、新文化”这些关键词；后来再看，才发现文章真正有力量的地方，不是把这些概念排列出来，而是从中国社会的实际状况一步一步推出道路选择。\footnote{毛泽东：《新民主主义论》，《毛泽东选集》第2卷，北京：人民出版社，1991年，第662--706页。}这让我第一次比较具体地理解了“实事求是”：它不是把现实说得很顺口，而是先承认现实的复杂，再从复杂里找方向。

过去我见到“实事求是”“群众路线”“独立自主”这些词，常常会自动把它们归到政治理论课的固定表达里，好像知道它们是什么意思就够了。可是这门课讲下来，我发现这些词如果只停留在解释层面，其实很容易变空。比如“独立自主”并不是关起门来什么都不学，而是在学习外来理论和经验时，不能把中国自己的历史条件丢掉。中国革命面对的不是已经完成资本主义发展的社会，而是一个被帝国主义、封建势力和官僚资本压迫的旧中国。用别人的现成道路解释这样的中国，必然会有不适合的地方。您当年提出把马克思主义基本原理同中国具体实际相结合，我现在理解，它首先是一种很艰难的判断能力，而不是一句后来总结出来的漂亮话。

课堂上讲到您的历史功绩，也讲到晚年探索中的曲折。我觉得这个部分反而让我更认真地看待历史人物。以前谈到您，我更容易想到开国领袖、革命胜利、社会主义制度建立这些宏大的叙述；但“历史功绩和晚年的悖论”这个专题提醒我，学习历史不能只选自己愿意看的那一面。新中国成立后，中国完成了从旧国家到新国家的巨大转变，建立了独立的工业体系、国防体系和基本制度框架，这些为后来的改革开放和现代化建设打下了基础。但与此同时，社会主义建设的探索也有失误和教训。把这两方面放在一起看，并不是为了抵消功绩，而是为了避免把历史讲得过于轻松。真正的历史理解，应该既能看见开创的艰难，也能承认探索中付出的代价。

我还想对您说，这门课让我对“人民”二字多了一点具体感觉。读《湖南农民运动考察报告》时，我最明显的感受是，您的判断不是从抽象概念里直接长出来的，而是从对社会现场的观察中来的。\footnote{毛泽东：《湖南农民运动考察报告》，《毛泽东选集》第1卷，北京：人民出版社，1991年，第12--44页。}过去我理解群众路线，常常把它当成“从群众中来，到群众中去”这句话本身；现在我会多想一步：如果没有真正看到普通人的处境、情绪和力量，就很难说自己理解了中国。人民不是文章里用来支撑论证的词，而是历史真正展开的地方。这个认识对我挺重要，因为它把宏大的历史叙述拉回到了具体的人身上。

把这些内容放到今天来看，我会想到新时代青年该怎样理解自己的责任。我们已经不处在救亡图存的战争年代，也不用以当年的方式回答“中国向何处去”。但是问题换了一种形式还在。今天的中国面对的是科技竞争、产业升级、文化传承、国家统一和国际环境变化等新的任务。对大学生来说，所谓“建设中国”不一定一开始就表现为很大的事业，它可能先体现在专业课有没有真正学进去，面对公共问题时有没有耐心了解事实，写一篇作业时有没有努力把自己的想法说清楚。以前我容易把课程、考试和未来工作分成几个互不相干的部分，这门课让我意识到，个人选择其实总会和时代处境发生关系，只是有时候自己没有认真想过。

当然，我也不想把自己的收获说得太满。学完这一学期，我还不能说自己已经真正理解了毛泽东思想，更不能说已经把中国道路讲清楚了。相反，我更清楚地感觉到，很多看似熟悉的词，如果不放回历史背景和现实问题中，就会变得很轻。比如“从实际出发”，说起来简单，真正做到却很难；“服务人民”听起来宏大，落实到日常选择时又常常很具体。也正因为如此，我觉得这门课的意义不只是让我掌握了几章内容，而是提醒我以后面对社会问题时，不要急着用简单结论把复杂现实盖住。

半个世纪后，我想对您说，您留给后人的不只是革命胜利的历史记忆，也包括一种不断追问现实、进入现实、改造现实的方法。它要求人既要有理想，又不能离开实际；既要相信方向，又要尊重规律；既要谈国家和民族，也要看见普通人的生活。对我这样的青年学生来说，这可能就是本学期最实在的收获：我还没有能力回答所有时代问题，但至少应该从自己的学习和生活开始，练习更认真地认识中国。

如果把这篇课程总结看作一次迟到五十年的对话，我现在能说的还很朴素：纪念您，不能只停留在熟悉的口号里，也不能把历史人物放成遥远的塑像。更重要的是，在今天的条件下继续学习怎样认识中国、怎样理解人民、怎样把个人成长放进国家发展中。未来我走向什么岗位，还不能完全确定，但我希望自己至少能够记住这一点：青年人的道路，不该只由个人得失来决定，也应当和这个国家正在回答的问题连在一起。这就是这门课结束时，我最想对您说的话。

\vspace{1em}
\noindent{\heiti 参考文献}

\noindent [1] 毛泽东：《新民主主义论》，《毛泽东选集》第2卷，北京：人民出版社，1991年。

\noindent [2] 毛泽东：《湖南农民运动考察报告》，《毛泽东选集》第1卷，北京：人民出版社，1991年。

\noindent [3] 毛泽东：《青年运动的方向》，《毛泽东选集》第2卷，北京：人民出版社，1991年。

}

\end{document}
